\section{Competir con los gigantes: los cuatro «no» y la ventana de tiempo}

La sección anterior trató la definición del producto; esta trata la competencia. Zhu Mingming (祝铭明) cita el resumen de Ma Yun (马云) sobre las cuatro oportunidades que tiene una empresa emergente frente a los gigantes: no lo ven, no lo entienden, no les interesa y no llegan a tiempo. En el caso de los lentes inteligentes, las grandes empresas sin duda lo ven y muy probablemente también lo entienden; la ventana de la empresa emergente proviene sobre todo de los dos o tres años que transcurren entre «no les interesa» y «no llegan a tiempo». La meta de Rokid no es competir en igualdad de condiciones con los gigantes, sino lograr una competencia entre pares, o al menos «sentarse a la mesa de los adultos».

\begin{figure}[H]
\centering
\includegraphics[width=0.96\textwidth]{startup-vs-giants-four-no.png}
\caption{Los cuatro «no» de la empresa emergente frente a los gigantes: no lo ven, no lo entienden, no les interesa, no llegan a tiempo. Diagrama conceptual propio, elaborado a partir de la conversación de 01:19:29--01:23:38.}
\end{figure}

\begin{knowledgebox}{Lectura de la figura: la ventana de Rokid no consiste en que los gigantes no lo sepan, sino en que todavía no se dedican de lleno}
En la figura, «no llegan a tiempo» es el estado que la empresa emergente realmente busca conseguir. Los lentes inteligentes son lo bastante visibles como para que las grandes empresas no los pasen por alto mucho tiempo; Rokid debe tomar suficiente ventaja en producto, usuarios, desarrolladores y cadena de suministro antes de que las grandes empresas entren de lleno.
\end{knowledgebox}

\begin{knowledgebox}{Términos clave: los cuatro «no»}
\noindent\begin{tabularx}{\textwidth}{p{0.16\textwidth}p{0.36\textwidth}X}
\toprule
Etapa & Situación del gigante & Ventana de la empresa emergente \\
\midrule
No lo ven & La oportunidad es demasiado marginal y la organización no la ha notado & Un equipo pequeño puede validar primero una necesidad nueva. \\
No lo entienden & Vieron el fenómeno, pero no pueden juzgar su valor & La empresa emergente demuestra su relevancia con producto y datos. \\
No les interesa & Lo entendieron, pero por ahora la escala no es suficiente & Es la ventana más probable hoy para los lentes inteligentes. \\
No llegan a tiempo & Cuando el gigante entra, el pionero ya acumuló experiencia y ecosistema & La empresa emergente debe buscar llegar a este estado. \\
\bottomrule
\end{tabularx}
\end{knowledgebox}

\subsection{Por qué las grandes empresas llegan un paso atrás}

Zhu Mingming considera que las grandes empresas suelen atreverse a ir detrás de los demás: entran con cadenas de suministro maduras, soluciones maduras y menor riesgo. Por ejemplo, los lentes sin pantalla son más fáciles de fabricar con la cadena de suministro existente; en cambio, los lentes inteligentes con pantalla, ligeros y all-in-one requieren explorar procesos de fabricación nuevos, y ahí la empresa emergente puede ponerse al frente. El problema es que ir adelante implica asumir, en nombre de toda la industria, los riesgos de los procesos de fabricación, de la capacidad de producción y de la educación del usuario.

\begin{warningbox}{Ser pionero no es solo una ventaja}
El pionero debe asumir procesos de fabricación no validados, el aumento gradual de la cadena de suministro, defectos de producto que se magnifican, la gestión de las expectativas del usuario y la presión financiera. Las grandes empresas, al entrar un paso después, pueden copiar la ruta ya madura; la empresa emergente tiene que convertir esa diferencia de tiempo en reputación, desarrolladores y experiencia de producto.
\end{warningbox}

\subsection{Competencia entre pares: no es competencia en igualdad de condiciones}

Esta subsección explica lo que Zhu Mingming llama «sentarse a la mesa de los adultos». Una empresa emergente no puede competir en igualdad de condiciones con Tencent (腾讯), Alibaba (阿里), ByteDance (字节), Xiaomi (小米), Apple (苹果) o Meta, porque el capital, el tráfico, la cadena de suministro y los canales no están en el mismo orden de magnitud. Lo que sí puede conseguir es una competencia entre pares: que, en una definición de producto concreta, los usuarios estén dispuestos a comparar Rokid con los gigantes dentro del mismo conjunto de opciones; que los canales estén dispuestos a venderlo en serio; que los desarrolladores estén dispuestos a adaptar sus aplicaciones, y que los socios de modelos y de contenido estén dispuestos a colaborar.

En la entrevista se mencionan varias señales para medirlo: uso diario promedio de más de dos horas, más de diez mil desarrolladores, más de tres mil desarrolladores empresariales, actividad semanal superior al setenta por ciento y una proporción elevada de recompra y recomendación. Estas cifras no deben tomarse como conclusiones auditadas externamente, sino como el tipo de evidencia que el ponente usa para mostrar que «se están formando atributos de plataforma». Lo que de verdad importa es que una empresa emergente no puede sentarse a la mesa solo con el entusiasmo de un evento de lanzamiento: tiene que demostrar tiempo de uso, recompra, desarrolladores y capacidad de entrega de la cadena de suministro.

\noindent\begin{tabularx}{\textwidth}{p{0.20\textwidth}p{0.32\textwidth}X}
\toprule
Indicador de competencia entre pares & Por qué importa & Limitaciones \\
\midrule
Tiempo de uso & Indica que el dispositivo no es una novedad pasajera & También hay que observar la estructura de tareas y la retención. \\
Recompra/recomendación & Indica que los primeros usuarios difunden el producto de boca en boca & La muestra puede estar sesgada hacia entusiastas de la tecnología; hay que ampliarla al usuario común. \\
Número de desarrolladores & Indica los primeros indicios de una oferta del ecosistema & Aún hay que validar la actividad y los ingresos de los desarrolladores. \\
Reabastecimiento de los canales & Indica que el producto se sostiene en ventas reales & La primera distribución no equivale a canales sólidos. \\
Capacidad de producción & Indica que el entusiasmo puede convertirse en entregas & Una capacidad alta también magnifica los problemas de calidad. \\
\bottomrule
\end{tabularx}

\subsection{Ecosistema y aliados}

Rokid no cree que pueda construir sola un ecosistema completo. Zhu Mingming menciona que Tencent tiene contenido, redes sociales y escenarios de uso, y que también pueden integrarse las capacidades de modelos como Doubao (豆包), Tongyi Qianwen (通义千问) y DeepSeek. Una empresa emergente de hardware debe consolidar la plataforma, la tecnología y la experiencia de producto, pero la construcción del ecosistema requiere más aliados. El juicio realista aquí es que, si los lentes inteligentes han de convertirse en un punto de acceso, no pueden ganar solo con especificaciones de hardware: también necesitan la participación conjunta de contenido, aplicaciones, desarrolladores, modelos y canales.

\begin{importantbox}{El ecosistema no lo construye una sola empresa}
Rokid puede asumir primero la exploración y el liderazgo, y consolidar la tecnología, la plataforma y la experiencia de producto; pero los socios de contenido, redes sociales, modelos, desarrolladores y escenarios de uso tienen que entrar en conjunto. Para un punto de acceso de IA portátil, si una empresa de hardware quiere hacer a la vez el hardware, el sistema operativo, los modelos, el contenido, las aplicaciones y los canales, el consumo de recursos supera rápidamente lo que una empresa emergente puede soportar.
\end{importantbox}

\subsection{Resumen de la sección}

Los lentes inteligentes son un campo de batalla al que los gigantes sin duda entrarán. La ventana de la empresa emergente está en asumir antes los riesgos, formar más rápido la experiencia de producto y su posición en el ecosistema, y ganarse el derecho a competir entre pares antes de que los gigantes entren de lleno.
